The event $ A = \{ \omega \in \Omega: \omega_1 + \omega_2 + \dots + \omega_{N} = N \} $  admits only
one elementary event, namely,
\[
\omega_{1}, \dots, \omega_{N} = 1 
.\] 
The probability of this event is 
\begin{align*}
    \P(A) = \P(\omega) &= q^{ \sum_{i=1}^{N} \frac{1+ \omega_{i} }{2} } (1-q)^{ \sum_{i=1}^{N}
\frac{1-\omega_{i} }{2} } \\
		       &= q^{ \sum_{i=1}^{N} \frac{1+1}{2} } (1-q)^{ \sum_{i=1}^{N} \frac{1-1}{2} }
		       \\
		       &= q^{ \sum_{i=1}^{N} 1 } (1-q)^{ \sum_{i=1}^{N} 0 } \\
		       &= q^{N} (1-q)^{0} = q^{N} 
\end{align*}

